\documentclass[a4paper]{article}
\usepackage{amsfonts}
\usepackage{amsmath}
\usepackage{amssymb}
\usepackage{graphicx}     

\begin{document}
  
\noindent \large Auteur: Anna Voronovska \\
\noindent \large Studierichting: Informatica\\
\noindent \large Oefeningenreeks 2B nr. 6.4\\

\medskip

\normalsize


$\operatorname{Bgtan} \dfrac{1}{2} + \operatorname{Bgtan} \dfrac{1}{5} + \operatorname{Bgtan} \dfrac{1}{8} = \dfrac{\pi}{4}$\\

Stel:\\ 

$\alpha = \operatorname{Bgtan} \dfrac{1}{2} \implies \tan \alpha = \dfrac{1}{2}$\\

$\beta = \operatorname{Bgtan} \dfrac{1}{5} \implies \tan \beta = \dfrac{1}{5}$\\

$\gamma = \operatorname{Bgtan} \dfrac{1}{8} \implies \tan \gamma = \dfrac{1}{8}$\\

$\tan(\alpha + \beta) = \dfrac{\dfrac{1}{2} + \dfrac{1}{5}}{1 - \dfrac{1}{2} \cdot \dfrac{1}{5}} = \dfrac{7}{9}$\\

$\tan(\alpha + \beta + \gamma) = \dfrac{\dfrac{7}{9} + \dfrac{1}{8}}{1 - \dfrac{7}{9} \cdot \dfrac{1}{8}} = 1$\\

$\tan(\alpha + \beta + \gamma) = 1$\\

$\tan\left(\dfrac{\pi}{4}\right) = 1$\\

$\alpha + \beta + \gamma = 1$\\

$\operatorname{Bgtan}(1) = \dfrac{\pi}{4}$\\


\begin{thebibliography}{999}
%\bibitem{a} Referentie
\end{thebibliography}
\end{document} 
